%=======================================================================
\section{Mathematical implementation}
\label{sec:implementation}
%=======================================================================

\subsection{Time integration}
\label{sec:integration}

The segment equation~\eqref{eq:ode} is nonlinear, because $\Tpcm(E')$ is
piecewise, but it is affine within each branch of
Eq.~\eqref{eq:enth-curve}. Over a time step $\Delta t$ the upstream water
temperature $T_j$ and the conductance $K_j$ are held at their values at the
beginning of the step, which are obtained by marching
Eqs.~\eqref{eq:Ui}, \eqref{eq:ntu}, \eqref{eq:K} and \eqref{eq:T-march} from
the inlet with the state frozen. The distinction between the branches governs
the choice of integrator.

\paragraph{Explicit stepping} On the latent plateau $\Tpcm=T_m$ does not depend
on $E'$, the segment behaves as though its heat capacity were infinite, and
Eq.~\eqref{eq:ode} reduces to $\mathrm{d}E'/\mathrm{d}t=K(T_j-T_m)$, for which
an explicit Euler step is exact. A model whose state is the melted area
occupies this regime permanently, which is why explicit stepping is adequate
for such models. On the sensible branches Eq.~\eqref{eq:ode} is a linear
relaxation with time constant $\tau=C/K$, and explicit Euler is stable only
for $\Delta t<2\tau$ and monotone only for $\Delta t<\tau$. The time grid
used here is logarithmic, which is the natural choice for a process
controlled by a growing conduction resistance, and its late steps exceed
$\tau$ by up to an order of magnitude \chk{ratio 7.7 at the last step of the
reference case, first segment; regenerate for the final grid}; an explicit
update diverges there.

\paragraph{Branchwise-exact scheme} Each branch is instead integrated in
closed form. With $T_j$ and $K$ fixed over the step,
\begin{equation}
  E'(t+\Delta t)=
  \begin{dcases}
    E'+K\left(T_j-T_m\right)\Delta t, & 0\le E'\le\Elat,\\[4pt]
    \Eeq+\left(E'-\Eeq\right)e^{-K\Delta t/C}, & \text{sensible branches},
  \end{dcases}
  \label{eq:exact}
\end{equation}
with the capacity and equilibrium enthalpy of the branch,
\begin{equation}
  \left(C,\,\Eeq\right)=
  \begin{cases}
    \left(C_s,\;C_s\left(T_j-T_m\right)\right), & E'<0,\\[3pt]
    \left(C_l,\;\Elat+C_l\left(T_j-T_m\right)\right), & E'>\Elat .
  \end{cases}
  \label{eq:Eeq}
\end{equation}
$\Eeq$ is the enthalpy at which the segment would be in equilibrium with the
water; the exponential approaches it monotonically and cannot overshoot it for
any $\Delta t$, so the scheme is unconditionally stable.

\paragraph{Branch crossings} If the state would leave its branch within the
step, the step is split at the crossing. With $E'_b\in\{0,\Elat\}$ the
boundary approached, the crossing time is
\begin{equation}
  t^{*}=
  \begin{dcases}
    \frac{E'_b-E'}{K\left(T_j-T_m\right)}, & \text{from the plateau},\\[6pt]
    \frac{C}{K}\ln\frac{E'-\Eeq}{E'_b-\Eeq}, & \text{from a sensible branch};
  \end{dcases}
  \label{eq:crossing}
\end{equation}
the state is advanced to $E'_b$, the remaining $\Delta t-t^{*}$ is integrated
in the adjacent branch, and the test is repeated (Table~\ref{tab:alg-segment}).
When the state is placed on a boundary it is displaced into the adjacent
branch by $\epsilon=10^{-9}\max(\Elat,1)$. Without that displacement a state
lying exactly on $\Elat$ is classified as two-phase by a closed-interval test,
the plateau branch returns $t^{*}=0$, and the segment remains pinned at the
boundary with zero superheat indefinitely.

\begin{table}[htbp]
\centering
\caption{Advance of one segment over one time step.}
\label{tab:alg-segment}
\small
\begin{tabular}{@{}r l@{}}
\toprule
1 & $t_\mathrm{left}\leftarrow\Delta t$;\quad $E'_{0}\leftarrow E'$\\
2 & \textbf{while} $t_\mathrm{left}>0$:\\
3 & \quad identify the branch of Eq.~\eqref{eq:enth-curve} containing $E'$\\
4 & \quad compute $t^{*}$ toward the boundary in the direction of $T_j-\Tpcm$, Eq.~\eqref{eq:crossing}\\
5 & \quad \textbf{if} $t^{*}\ge t_\mathrm{left}$: advance by Eq.~\eqref{eq:exact} over $t_\mathrm{left}$; \textbf{stop}\\
6 & \quad \textbf{else}: $E'\leftarrow E'_b\pm\epsilon$;\quad $t_\mathrm{left}\leftarrow t_\mathrm{left}-t^{*}$\\
7 & \textbf{return} $E'$ and $q'=(E'-E'_{0})/\Delta t$\\
\bottomrule
\end{tabular}
\end{table}

\paragraph{Energy closure} The heat rate passed to the water is not evaluated
from Eq.~\eqref{eq:K} independently of the state update; it is the
step-averaged value returned by the integrator,
$q'_j=(E'^{\,n+1}_j-E'^{\,n}_j)/\Delta t$, and the water temperature is reduced
by exactly that amount through Eq.~\eqref{eq:T-march}. The energy gained by
the PCM therefore equals that lost by the water to round-off, and the second
equalities in Eqs.~\eqref{eq:Ech} and~\eqref{eq:Edc} are exact. This is a
consistency property of the implementation and not evidence of its
correctness. The integrator itself was verified against an explicit
integration of Eq.~\eqref{eq:ode} with $2\times10^{5}$ substeps over a full
charge and discharge, including branch crossings, to a relative difference of
$3\times10^{-13}$ \chk{re-run on v0.12 before quoting}.

%-----------------------------------------------------------------------
\subsection{Cyclic steady state and the sizing problem}
\label{sec:css}
%-----------------------------------------------------------------------

A store operated on a daily schedule does not begin each charge from a fully
solid state. The charge--discharge pair is therefore repeated, each charge
starting from the state left by the preceding discharge, until the store
returns to its own initial state,
\begin{equation}
  \max_{j}\left|E'_j\big|_{\text{end of cycle }n}-E'_j\big|_{\text{start of cycle }n}\right|
  <\epsilon_\mathrm{CSS},
  \label{eq:css-test}
\end{equation}
with $\epsilon_\mathrm{CSS}=\SI{e-9}{\joule\per\meter}$. The first cycle
starts from a fully solid store at the melting temperature of each layer,
$E'_j=0$.

At cyclic steady state the enthalpy change of the store over a complete cycle
is zero by definition, and under H6 the store exchanges heat only with the
water. Hence
\begin{equation}
  \eta_\mathrm{storage}\equiv\frac{\Delta E_\mathrm{dc}}{\Delta E_\mathrm{ch}}=1
  \label{eq:css-unity}
\end{equation}
identically, for every well count and every pair of flow rates. This is
energy conservation applied to a closed cycle in a lossless model, not a
result about the store, but it determines how the design problem may be posed.
A single-cycle model, such as that of Ref.~\cite{KoIHTC2026}, sizes the field
by requiring that it absorb $\Delta E_\mathrm{out,HP}$ within $t_\mathrm{ch}$
and solves for $\Nw$ (and, separately, for the discharging flow). At CSS the
absorbed and released energies are the same number, so the charging and
discharging requirements collapse into a single equation. That equation
fixes the ratio of the two flow rates and places no constraint on $\Nw$: the
sizing problem has no root. Introducing a loss mechanism would not remove the
degeneracy, since at CSS the charge would then exceed the discharge by exactly
the loss and the one remaining equation would still determine only the flow
ratio.

The problem is therefore posed as a simulation. The field is specified by the
latent inventory alone,
\begin{equation}
  \Nw=\frac{\Delta E_\mathrm{out,HP}}{\rho_s\,\Vwell\,\hm},
  \label{eq:Nwell}
\end{equation}
which ignores the sensible capacity of the store and is to that extent
conservative; the two flow rates are fixed by their glides,
Eq.~\eqref{eq:mdot}; the field is marched to CSS; and the relative deviation
of the delivered energy from the requirement,
\begin{equation}
  \mathcal{D}=\frac{\Delta E_\mathrm{dc}}{\Delta E_\mathrm{in,ORC}}-1,
  \label{eq:deviation}
\end{equation}
is reported as a result. Because $\md_\mathrm{dc}$ is fixed by the specified
glide, the delivered energy over a fixed window is, for constant $c_p$,
\begin{equation}
  \frac{\Delta E_\mathrm{dc}}{\Delta E_\mathrm{in,ORC}}
  \simeq\frac{\overline{T}_\mathrm{2d}-T_\mathrm{3d}}{\Delta T_\mathrm{3D,2D}},
  \label{eq:glide-shortfall}
\end{equation}
so $\mathcal{D}$ is the shortfall (or excess) of the glide the field actually
achieves relative to the one assumed \chk{v0.13 reference case: 1.04458
against 1.04413; the residual difference reflects the variation of $c_p$ with
temperature}. Two consequences
follow. First, $\mathcal{D}$ does not enter RTE: the charging flow is pinned
in the same way, so a field that delivers \SI{1}{\percent} more than the
target also absorbs \SI{1}{\percent} more, and $\mathcal{D}$ measures plant
size rather than performance. Second, $\mathcal{D}$ is not to be driven to
zero by adjusting a design parameter, since doing so would convert the one
independent output of the procedure back into an imposed constraint.

The residual melted fraction at the end of discharge,
$\bar\varepsilon_\mathrm{dc}$, is used as a diagnostic of the limiting
mechanism. If it is close to zero the store empties within the window and the
field is limited by its inventory; if it is appreciably above zero, and
particularly if it rises with $\Nw$, the field is limited by the rate of heat
transfer, and the effective levers are the flow rate, the window length and
the exchanger conductance rather than the PCM inventory.

%-----------------------------------------------------------------------
\subsection{Verification of the exergy audit}
\label{sec:ex-verification}
%-----------------------------------------------------------------------

The component destructions of Eqs.~\eqref{eq:sgen-hp}--\eqref{eq:sgen-orc} are
computed from the entropy generated inside each component. For each machine
the sum is compared with a boundary balance constructed from stream states
alone,
\begin{align}
  \sum_{j\in\mathrm{HTHP}}I_j&=W_\mathrm{C}+\left(\Ex_\mathrm{geo}-\Ex_\mathrm{reinj}\right)
    -\Exd_\mathrm{w}^\mathrm{tot}\!\left(T_\mathrm{3c},T_\mathrm{2c}\right)t_\mathrm{ch}
    +I_\mathrm{mot},
  \label{eq:gate-hp}\\
  \sum_{j\in\mathrm{ORC}}I_j&=\Exd_\mathrm{w}^\mathrm{tot}\!\left(T_\mathrm{2d},T_\mathrm{3d}\right)t_\mathrm{dc}
    -W_\mathrm{T}+I_\mathrm{gen},
  \label{eq:gate-orc}
\end{align}
and the two sides agree to $10^{-15}$ in relative terms for the reference
case. These two checks are independent and can fail: they detected the
inconsistency in the water flow rates discussed after Eq.~\eqref{eq:mdot}.
The global balance, Eq.~\eqref{eq:ex-balance}, is by contrast an identity
under the present construction. Because $T_\mathrm{4c}=T_\mathrm{3c}$ and the
discharge terminal temperatures are shared with the ORC evaporator,
Eq.~\eqref{eq:I-BB} cancels the water terms of the two surface exchangers
exactly, and the global residual is zero by construction rather than by
verification. $I_\mathrm{BB}$ is therefore correct but not independently
verified until it is evaluated from $s(E')$.

Two approximations that the audit carried in earlier versions have since been
removed. It was evaluated on the target duties of Section~\ref{sec:budget}
rather than on the energy the field actually moves, and the pumping work was
excluded altogether; the round-trip efficiency used the realised duty and
charged the pumping, so the two quantities did not share a control volume and
$\psi$ could not properly be compared with RTE. Both are now on the
cyclic-steady-state basis. Because the storage efficiency is identically unity
at CSS and $\lambda$ is zero, the two half-cycles scale together and the
rebase is a single factor, \num{1.0446} in the reference case; the pumping
work enters as an input and its dissipation in the tubes as a destruction,
\SI{703}{\kilo\watt\hour} per cycle, or \SI{2.7}{\percent} of the total.

One approximation remains: the audit uses the nominal water terminal
temperatures rather than the flow-averaged outlets
$\overline{T}_\mathrm{2c}$ and $\overline{T}_\mathrm{2d}$ that the field
actually returns. The resulting mismatch in water exergy is reported by the
audit as a separate diagnostic rather than absorbed into any component.

%-----------------------------------------------------------------------
\subsection{Solution procedure}
\label{sec:procedure}
%-----------------------------------------------------------------------

The procedure is summarised in Fig.~\ref{fig:procedure}. It is organised in
three nested levels. At the plant level the ORC and the HTHP are evaluated
once (Section~\ref{sec:cycles}), and their outputs to the rest of the
calculation are $\etaORC$, $\COP$ and the secondary-fluid temperatures of
Table~\ref{tab:closure}. At the field level the target output fixes the
energy budget, Eqs.~\eqref{eq:E-orc}--\eqref{eq:W-in} with $\lambda=0$, the
well count, Eq.~\eqref{eq:Nwell}, and the flow rates, Eq.~\eqref{eq:mdot}.
At the borehole level the charging and discharging marches are repeated until
Eq.~\eqref{eq:css-test} is satisfied. No quantity is obtained by root finding;
the loop to CSS is the only iteration.

Because $T_\mathrm{3e}$ depends on the realised outlet $\overline{T}_\mathrm{2d}$,
which is available only after the march, the whole procedure is executed
twice: first with the provisional outlet $T_\mathrm{2d}^{(0)}$, then with the
flow-averaged outlet from the first pass. The ORC, the energy budget, $\Nw$
and both flow rates are re-evaluated in the second pass. A single correction
suffices because the outlet is controlled by the store rather than by the
plant; the change in $\overline{T}_\mathrm{2d}$ between the two passes is
reported as a check \chk{below \SI{e-4}{\kelvin} in the v0.13 reference case}.

\begin{figure}[htbp]
\centering
\begin{tikzpicture}[
  font=\footnotesize, node distance=2.6mm,
  blk/.style ={draw, rounded corners=2pt, inner sep=4pt, align=left,
               text width=0.80\linewidth},
  inp/.style ={blk, fill=black!4},
  mar/.style ={blk, line width=0.9pt, fill=black!2},
  tst/.style ={draw, dashed, rounded corners=2pt, inner sep=4pt, align=center,
               text width=0.62\linewidth},
  fl/.style  ={-{Latex[length=2mm]}, thick},
  lab/.style ={font=\scriptsize, inner sep=1.5pt}]

\node[inp] (in) {\textbf{Inputs.} $\dot W_\mathrm{el,out}$, $t_\mathrm{ch}$,
  $t_\mathrm{dc}$; approaches and glides of Table~\ref{tab:closure};
  $T_\mathrm{source}$, $T_\mathrm{sink}$; isentropic, generator and motor efficiencies; PCM
  properties, $T_m^\mathrm{top}$, $\Nlay$, $\Delta T_\mathrm{casc}$; well and
  tube geometry; working fluids; $n_s$, $n_t$ time levels,
  $\epsilon_\mathrm{CSS}$};

\node[blk, below=of in] (cyc) {\textbf{1.\ Thermodynamic cycles}
  \hfill\emph{plant level}\\
  cascade, Eq.~\eqref{eq:cascade}; inlets $T_\mathrm{4c}$, $T_\mathrm{3d}$,
  Eq.~\eqref{eq:inlets}; $T_\mathrm{2d}$ (provisional, then realised)\\
  ORC: Eqs.~\eqref{eq:orc-p}--\eqref{eq:T3e}, $T_\mathrm{3e}$ by bisection on
  the pinch $\Rightarrow\etaORC$\\
  HTHP: Eqs.~\eqref{eq:hp-p}--\eqref{eq:hp-T}, suction states by bisection,
  $x_\mathrm{8h}$ by fixed point $\Rightarrow\COP$};

\node[blk, below=of cyc] (bud) {\textbf{2.\ Energy budget and hardware}
  \hfill\emph{field level, no root finding}\\
  $\dot Q_\mathrm{in,ORC}$, $\Delta E_\mathrm{in,ORC}$, $\Delta
  E_\mathrm{out,HP}$ ($\lambda=0$), Eqs.~\eqref{eq:E-orc}--\eqref{eq:W-in};
  $\Nw$, Eq.~\eqref{eq:Nwell}; $\md_\mathrm{ch}$, $\md_\mathrm{dc}$,
  Eq.~\eqref{eq:mdot}};

\node[mar, below=of bud] (ch) {\textbf{3.\ Charge}\hfill\emph{borehole level}\\
  inlet $T_\mathrm{4c}$ at $s=0$, flow $\md_\mathrm{ch}$; \textbf{for} each
  time level $t_k$ (logarithmic):\\
  \hspace*{3mm}\textbf{for} $j=1,\ldots,n_s$: state $\to\Amelt,\Tpcm$;
  shell, Eq.~\eqref{eq:shell-choice} or $\delta_\mathrm{eq}$;
  $h_i$, $\etao$, $\Ui$; $\NTU_j$, $K_j$;\\
  \hspace*{9mm}advance $E'_j$ (Table~\ref{tab:alg-segment});
  $T_{j+1}$ from Eq.~\eqref{eq:T-march}};

\node[mar, below=of ch] (dc) {\textbf{4.\ Discharge}\\
  mirror cascade and state; inlet $T_\mathrm{3d}$ at $s=\Ltube$, flow
  $\md_\mathrm{dc}$; march as in 3 with the freezing shell};

\node[tst, below=of dc] (tst) {state returned to itself, Eq.~\eqref{eq:css-test}?};

\node[blk, below=5mm of tst] (cor) {\textbf{5.\ Correction pass}
  \hfill\emph{executed once}\\
  repeat 1--4 with $T_\mathrm{2d}=\overline{T}_\mathrm{2d}$ from the first
  pass; report the change in $\overline{T}_\mathrm{2d}$};

\node[inp, below=of cor] (out) {\textbf{6.\ Results at CSS.}
  $\Delta E_\mathrm{ch}$, $\Delta E_\mathrm{dc}$, deviation $\mathcal D$,
  Eq.~\eqref{eq:deviation}; $\overline{T}_\mathrm{2c}$,
  $\overline{T}_\mathrm{2d}$; $\bar\varepsilon_\mathrm{ch}$,
  $\bar\varepsilon_\mathrm{dc}$; pumping, Eq.~\eqref{eq:Wf}; $\eta_T$,
  $\eta_{T,\mathrm{eff}}$, RTE, $\varepsilon_\mathrm{RTE}$,
  $\Delta E_\mathrm{therm}$, $\rho_E$; $UA$ of the plant exchangers,
  Eq.~\eqref{eq:UA}; exergy audit, $\psi$, $\md_\mathrm{geo}$,
  $N_\mathrm{geo}$, Section~\ref{sec:exergy}; profiles of $T$, $\Tpcm$,
  $\eloc$, $\Ui$, $\NTU$};

\foreach \a/\b in {in/cyc,cyc/bud,bud/ch,ch/dc,dc/tst,cor/out}
  \draw[fl] (\a) -- (\b);
\draw[fl] (tst) -- node[lab,right] {yes} (cor);
\coordinate (L) at ($(ch.west)+(-7mm,0)$);
\draw[fl] (tst.west) -- (L |- tst.west)
  node[lab,midway,above] {no: next cycle} -- (L) -- (ch.west);
\end{tikzpicture}
\caption{Solution procedure. Blocks 1--2 are evaluated once per pass and
produce a fully specified field; blocks 3--4 are repeated until cyclic steady
state; block 5 is a second evaluation, not a loop.}
\label{fig:procedure}
\end{figure}

%-----------------------------------------------------------------------
\subsection{Simulation conditions}
\label{sec:conditions}
%-----------------------------------------------------------------------

The reference case is summarised in Table~\ref{tab:conditions}. The PCM
properties are consistent with the trimodal material of Saher \textit{et al.}\
\cite{Saher2024}: a top-layer melting temperature of \SI{150}{\celsius} and a
latent heat of \SI{380}{\kilo\joule\per\kilogram}; in the absence of reported
data, densities of \num{1550} and \SI{1450}{\kilogram\per\meter\cubed},
specific heats of \num{1280} and \SI{1800}{\joule\per\kilogram\per\kelvin}
and conductivities of \num{0.60} and \SI{0.45}{\watt\per\meter\per\kelvin}
are assumed for the solid and liquid, respectively. \chk{The cascade assumes a
family of PCMs of identical properties with melting points spanning
\SIrange{101}{150}{\celsius}; state this explicitly and discuss whether such a
family exists.} The wells are \SI{5000}{ft} (\SI{1524}{\meter}) deep with a
\SI{7}{in} (\SI{177.8}{\milli\meter}) bore, each containing two hairpins of
1\,$\tfrac14$\,in Schedule~80 tube with 24 fins per leg, which leaves
$\Vwell=\SI{27.68}{\meter\cubed}$ of PCM per well. The cell radius is
$r_\mathrm{cell}=\SI{44.45}{\milli\meter}$, so
$\delta_\mathrm{merge}=\SI{23.37}{\milli\meter}$, $\beta=2.109$,
$S=0.3467$ and $\delta_\mathrm{eq}=\SI{8.74}{\milli\meter}$.

\begin{table}[htbp]
\centering
\caption{Reference case.}
\label{tab:conditions}
\small
\resizebox{\linewidth}{!}{%
\begin{tabular}{@{}l l r@{}}
\toprule
 & quantity & value \\
\midrule
\multirow{4}{*}{Duty}
 & net electrical output $\dot W_\mathrm{el,out}$ & \SI{1}{\mega\watt} \\
 & charging / discharging window $t_\mathrm{ch}$, $t_\mathrm{dc}$ & \SI{10}{\hour} / \SI{10}{\hour} \\
 & HTHP source $T_\mathrm{source}$ / ORC sink $T_\mathrm{sink}$ & \SI{60}{\celsius} / \SI{20}{\celsius} \\
 & working fluid (ORC and HTHP) & cyclopentane \\
\midrule
\multirow{4}{*}{Machines}
 & isentropic: turbines $\eta_{s,t}$ / ORC pumps $\eta_{s,p}$ / compressors $\eta_{s,c}$ & 0.85 / 0.85 / 0.85 \\
 & generator $\eta_\mathrm{gen}$ / motor $\eta_\mathrm{mot}$ & 0.95 / 0.95 \\
 & ORC intermediate pressure & arithmetic mean \\
 & HTHP intermediate pressure & geometric mean \\
\midrule
\multirow{6}{*}{Approaches [K]}
 & $\Delta T_\mathrm{pinch,ORC}$ / $\Delta T_\mathrm{2D,3E}$ & 5 / 12 \\
 & $\Delta T_\mathrm{E,sink}$ / $\Delta T_\mathrm{sink,gl}$ & 5 / 0 \\
 & $\Delta T_\mathrm{2H,3C}$ / $\Delta T_\mathrm{sub}$ & 10 / 2 \\
 & $\Delta T_\mathrm{3A,4A}$ / $\Delta T_\mathrm{pinch,HPE}$ & 10 / 5 \\
 & $\Delta T_\mathrm{4C,M}$ / $\Delta T_\mathrm{M,1D}$ & 10 / 10 \\
 & glides $\Delta T_\mathrm{3C,2C}$ / $\Delta T_\mathrm{3D,2D}$ / $\Delta T_\mathrm{casc}$ & 55 / 55 / 55 \\
\midrule
\multirow{4}{*}{PCM}
 & $T_m^\mathrm{top}$ / $T_m^\mathrm{bot}$ / $\Nlay$ & \SI{150}{\celsius} / \SI{101.1}{\celsius} / 9 \\
 & $\hm$ & \SI{380}{\kilo\joule\per\kilogram} \\
 & $\rho_s$ / $\rho_l$; $c_{p,s}$ / $c_{p,l}$ & \SI{1550}{} / \SI{1450}{\kilogram\per\meter\cubed}; 1280 / \SI{1800}{\joule\per\kilogram\per\kelvin} \\
 & $k_s$ / $k_l$ & 0.60 / \SI{0.45}{\watt\per\meter\per\kelvin} \\
\midrule
\multirow{5}{*}{Geometry}
 & $\Lwell$ / $D_\mathrm{well}$ & \SI{1524}{\meter} / \SI{177.8}{\milli\meter} \\
 & hairpins per well $n_t$ & 2 \\
 & $r_i$ / $r_e$ / $k_\mathrm{wall}$ & \SI{16.23}{\milli\meter} / \SI{21.08}{\milli\meter} / \SI{45}{\watt\per\meter\per\kelvin} \\
 & $n_f$ / $L_f$ / $t_f$ & 24 / \SI{7.5}{\milli\meter} / \SI{1.5}{\milli\meter} \\
 & fouling $R''_{f,i}$ & \SI{e-3}{\meter\squared\kelvin\per\watt} \\
\midrule
\multirow{4}{*}{Geothermal source}
 & yield per producing well $\md_\mathrm{geo,well}$ & \SI{12.5}{\kilo\gram\per\second} \chk{provisional} \\
 & minimum reinjection temperature $T_\mathrm{reinj,min}$ & \SI{25}{\celsius} \chk{placeholder} \\
 & dead state $T_0$ & $T_\mathrm{sink}=\SI{20}{\celsius}$ \\
 & exergy convention & resource \\
\midrule
\multirow{3}{*}{Loop / numerics}
 & water pressure $p_\mathrm{w}$ & \SI{1}{\mega\pascal} \\
 & segments $n_s$ / time levels per half-cycle & 100 / 40 \chk{use 99 or 108} \\
 & first time level / CSS tolerance & \SI{1}{\second} / \SI{e-9}{\joule\per\meter} \\
\bottomrule
\end{tabular}}
\end{table}

The time levels are distributed logarithmically between \SI{1}{\second} and
the end of each window, $t_k=t_\mathrm{w}^{(k-1)/(n_t-1)}$ with $t_\mathrm{w}$
in seconds, so the last step of a \SI{10}{\hour} window is
\SI{8.5e3}{\second}. \chk{The warning raised by the code for $n_s=100$ and
$\Nlay=9$: a layer boundary falls inside a segment, which then carries a
melting temperature in error by up to one layer width (\SI{6.1}{\kelvin}).
Rerun with $n_s=99$ or $108$ for the paper, and report a grid-convergence
check in $n_s$ and $n_t$.}

\paragraph{Validity of H7 in the reference case} At the charging flow the
mean velocity in the tube is about \SI{1.25}{\meter\per\second}, so the water
transits the \SI{3048}{\meter} hairpin in about \SI{2.4e3}{\second}, or
\SI{7}{\percent} of the window. The heat capacity of the water in the tube is
\SI{26}{\percent} of the sensible capacity $C_l$ of the PCM it serves, but a
\SI{10}{\kelvin} swing of the water holds only \SI{1.2}{\percent} of the
latent inventory, so the unresolved start-up transient displaces little
energy although it occupies a noticeable fraction of the window.
\chk{Earlier drafts quote \SI{2264}{\second} and \SI{1.35}{\meter\per\second};
v0.11 and v0.12 give \SI{2438}{\second} and \SI{1.25}{\meter\per\second}.}

The yield of a geothermal producer is taken provisionally as
\SI{12.5}{\kilo\gram\per\second}, the value for a \SI{4500}{\meter},
\SI{206}{\celsius} repurposed well, which is deeper and hotter than the source
considered here; reported values for repurposed oil and gas wells range from
about \SIrange{1}{12.5}{\kilo\gram\per\second}. $N_\mathrm{geo}$ is inversely
proportional to this value and should be read as a scaling result.
\chk{add references for the per-well yield range and the \SI{4500}{\meter}
case study.}

The equations were implemented in Python with CoolProp \cite{CoolProp} for
fluid properties. \todo{Python/CoolProp versions; code availability
statement.}
